\section{Introduction}
\label{sec:intro}

Since 2012, global correspondent banking relationships have declined by over 25\%, severing payment corridors between dozens of countries \citep{Rice2020_correspondent}. This retreat, termed ``de-risking,'' accelerated after 2016 as banks withdrew from jurisdictions deemed high-risk for money laundering and terrorism financing \citep{Borchert2024_broken_relationships}. The consequences extend beyond financial access: when payment links break, the products requiring the most coordination and trust suffer the most.

This paper asks two questions. First, do payment-friction shocks differentially affect trade in complex products, and is the effect symmetric: do friction \emph{reductions} raise complex trade just as friction \emph{increases} reduce it? Second, where does stablecoin adoption (dollar-denominated tokens transacted on blockchain networks) offer the greatest opportunity to substitute for missing payment infrastructure?

The mechanism is straightforward. Complex products such as electronics, machinery, and pharmaceuticals require relationship-specific investment, multi-party coordination, and extended payment terms \citep{Manova2013_credit_constraints, AntrasFolley2015_trade_finance}. These products depend on dense networks of financial intermediaries to manage letters of credit, escrow, and foreign exchange conversion. When correspondent-banking relationships collapse, the marginal cost increase falls disproportionately on complex goods. The mirror image is also true. When jurisdictions harmonize KYC and AML standards, the per-corridor compliance burden falls disproportionately for products that pass through more compliance gates per shipment. The same friction channel that hurts complex trade in one direction helps it in the other.

We estimate a structural gravity model using Poisson PML (PPML) on bilateral trade flows among 226 countries in 94 product categories from 2015 to 2023 \citep{SantosSilva2006_ppml}. Identification rests on three reduced-form natural experiments. The first is the FATF greylist, a country-level designation issued after a multi-year compliance review by the Financial Action Task Force; greylisting raises correspondent-banking and KYC costs and is unrelated to bilateral goods trade. The second is the EU Anti-Money-Laundering Directive (AMLD4), effective in 2017, which harmonized AML and KYC standards across EU member states. The third is the FDI-derisked corridor measure: bilateral corridors where foreign direct investment positions dropped by more than 50\% relative to the 2015--2017 base period.

The three experiments deliver a tight causal triangle. In our pair-fixed-effects panel covering 2018--2023 and 2.82 million observations, FATF greylisting reduces trade in complex products differentially ($\hat{\gamma}_{\text{FATF}} = -0.108$, $p = 0.098$), AMLD harmonization \emph{raises} it ($\hat{\gamma}_{\text{AMLD}} = +0.141$, $p = 0.040$), and FDI de-risking reduces it ($\hat{\gamma}_{\text{FDI}} = -0.135$, $p = 0.043$). Yearly cross-sections from 2018 to 2023 confirm the pattern: FATF $\times$ PCI is negative in 6/6 years, AMLD $\times$ PCI is positive and significant at the 1\% level in 6/6 years, and FDI-derisked $\times$ PCI is negative and significant at the 1\% level in 6/6 years. A triple horse race that estimates all three interactions jointly with pair and product--year fixed effects preserves all three signs; AMLD ($p=0.050$) and FDI ($p=0.075$) remain significant while FATF retains its sign but loses precision in the joint specification ($\hat{\gamma}_{\text{FATF}} = -0.076$, $p = 0.166$), consistent with FATF and FDI capturing overlapping slices of the same payment-friction increase. Diplomatic disagreement, included as a placebo geopolitical channel, enters with a positive sign in the yearly kitchen-sink cross-sections (and is statistically null in the panel kitchen-sink), ruling out the interpretation that the complexity penalty is a geopolitical artifact.

The economic magnitudes are substantial. A one-standard-deviation more complex product loses roughly 11--30\% of its trade flow when the corridor enters the FATF greylist, and gains 10--15\% in EU corridors after AMLD. The mechanism is symmetric: payment-friction increases hurt complex trade, and payment-friction reductions help it.

Building on these estimates, we construct a Stablecoin Opportunity Score (SOS) that ranks countries by potential gains from alternative payment infrastructure. The SOS is a composite weighted-evidence aggregator across two distinct natural experiments: the FATF (clean, binary, exogenous) and FDI (broad, continuous, persistent) channels. Because the two single-instrument rankings are themselves weakly correlated ($\rho = 0.249$), the composite is not a robustness check on the rank order; it is the headline because it integrates evidence from both shocks. EU corridors enter with a negative AMLD adjustment: regulatory clarity is a partial substitute for stablecoins, so countries already covered by AMLD harmonization are down-weighted. The composite top five are Haiti, Fiji, Syria, Lebanon, and Nicaragua, small economies with high greylisting and de-risking exposure and limited regulatory harmonization to substitute. We classify countries into a feasibility quadrant based on capital-account openness and existing crypto on-ramp infrastructure, identifying 60 ``actionable'' economies where stablecoin deployment faces the fewest barriers; within that quadrant Panama, Trinidad and Tobago, the United Arab Emirates, Bahrain, and Georgia top the SOS ranking.

This paper makes three contributions. First, two distinct natural experiments---FATF greylisting (friction up) and EU AMLD harmonization (friction down)---establish a symmetric mechanism linking payment frictions to complex-product trade. The symmetry is the cleanest causal-mechanism design available in this setting: FATF is the natural-experiment plaintiff, and AMLD is the natural-experiment defendant. The two extend the financial-frictions--trade literature \citep{Manova2013_credit_constraints} by showing that payment-specific (not credit-specific) frictions interact with product complexity, and they extend the economic-complexity literature \citep{Hidalgo2009_building_blocks, Hausmann2014_atlas} by identifying payment infrastructure as a binding constraint on complex-product trade. Second, we provide the first product-level mapping of stablecoin opportunity, bridging the emerging literature on crypto capital flows \citep{GrafvonLuckner2023_capital_flows, Auer2025_defiying_gravity} with the structural gravity framework used to evaluate payment-system interlinking \citep{FerrarMinesso2026_payment_links}. Third, we develop a feasibility-weighted policy tool that identifies where stablecoin deployment is both high-impact and implementable, with the AMLD adjustment encoding that regulatory harmonization and stablecoin payment rails are partial substitutes.

The paper proceeds as follows. Section~\ref{sec:literature} reviews the related literature. Section~\ref{sec:data} describes the data. Section~\ref{sec:strategy} presents the empirical strategy. Section~\ref{sec:results} reports the main results. Section~\ref{sec:robustness} presents robustness checks and falsification tests. Section~\ref{sec:sos} constructs and validates the Stablecoin Opportunity Score. Section~\ref{sec:conclusion} concludes.
